% !TeX root = ../main.tex
% Lista de acrónimos (paquete acronym). En el texto:
%   \ac{ODS}   -> primera vez "Objetivos de Desarrollo Sostenible (ODS)", luego "ODS"
%   \acs{ODS}  -> siempre la sigla     \acl{ODS} -> siempre el nombre largo
%   \acp{...}  -> plural (añade una "s"; si el plural es irregular usa \acs)
% En títulos y pies de figura escribe la sigla como texto normal: si usas
% \ac o \acs, la "primera aparición" ocurre en el índice.
% Mantén la lista en orden alfabético. El argumento opcional de
% \begin{acronym} es la sigla más larga (para alinear la columna).
\cleardoublepage
\chapter*{Lista de acrónimos}
\addcontentsline{toc}{chapter}{Lista de acrónimos}
\begin{acronym}[ODS]
  \acro{CI}{integración continua (\textit{Continuous Integration})}
  \acro{ODS}{Objetivos de Desarrollo Sostenible}
  \acro{SI}{Sistema Internacional de Unidades}
  \acro{TFG}{Trabajo Fin de Grado}
  \acro{TFM}{Trabajo Fin de Máster}
\end{acronym}
